% Figure captions for the main manuscript figures
\textbf{ERP grand-average ROI.} Grand-average posterior-region ERP waveforms for target and non-target trials, averaged across PO7, PO8, PO3, PO4, O1, and O2. Shaded bands denote the 95\% confidence interval across epochs, and the highlighted 300--500 ms region corresponds to the expected target-related P300 interval.

\textbf{ERP difference wave ROI.} Target-minus-non-target difference waveform for the posterior ROI. The annotated peak within 250--600 ms summarizes the dominant target-related positivity in the RSVP task.

\textbf{Topographic difference map (300--500 ms).} Sparse posterior scalp topography of the target-minus-non-target ERP amplitude averaged over 300--500 ms. Because the montage contains only eight electrodes, the map should be interpreted as a coarse posterior topographic summary rather than a dense spatial reconstruction.

\textbf{Channel-time importance heatmap.} Model-derived feature importance aggregated by channel and ERP time window. This view highlights whether discriminative information clusters in posterior electrodes and in the canonical P300 interval.

\textbf{Precision-recall curve.} Precision-recall curve for the best pooled classifier. The dashed horizontal line indicates the positive-class prevalence, which provides the relevant baseline under class imbalance.

\textbf{Latency breakdown.} Average latency of the main processing stages, including preprocessing, epoching, feature extraction, and single-sample inference, illustrating the computational cost of the lightweight ERP-based pipeline.
